
\section{Related Work}
\paragraph{Test time scaling and latent reasoning.}
Test-time compute scaling is a critical frontier for enhancing LLM reasoning~\citep{zhang2025survey}.
Conventional approaches primarily rely on explicit reasoning, where models generate natural language intermediate steps such as Chain-of-Thought prompting~\citep{wei2022chain} to solve complex tasks. 
 % Recent advancements demonstrate that scaling the length of CoT sequences during inference significantly improves performance on reasoning benchmarks.
 Further gains have been achieved through search-based strategies, including majority voting~\citep{wang2022selfconsistency}, Tree-of-Thoughts (ToT)~\citep{tot}, and Monte Carlo Tree Search (MCTS)~\cite{yang2022chain}, which allow the exploration of multiple reasoning paths. However, these methods remain inherently constrained by the bandwidth and efficiency of natural language generation.
Inspired by the human tendency to reason through internal steps rather than producing immediate verbal outputs~\citep{zelikman2024quiet}, recent work has pursued latent reasoning, which shifts computation from discrete tokens into latent representations~\citep{zhu2025survey}.  This approach is more computationally efficient and better suited for abstract reasoning.
% One line of research introduces non-interpretable tokens to serve as “internal monologues,” enabling the model to process intermediate states without expanding the observable output. 
Coconut~\citep{hao2024training} employing continuous thought tokens derived from previous hidden states is a typical technique. 
 Approaches such as SIM-CoT~\citep{wei2025sim} and others~\citep{mohtashami2023cotformer,shen2025codi,cheng2024compressed,zhang2025lightthinker} share similar ideas. Nevertheless, these approaches essentially remain forms of test-time scaling along the sequential dimension.
\paragraph{Recurrent Transformers and LoopLM.}
Unlike methods that scale sequence length, a newer paradigm of latent reasoning focuses on scaling model depth through recurrence and parameter sharing.  Universal Transformer~\citep{dehghani2018universal} pioneered dynamic recurrence across layers, establishing depth-adaptive computation as an alternative to fixed-depth transformers.
Subsequent research~\citep{giannou2023looped,yang2023looped,fan2024looped} have investigated their potential benefits through theoretical and small-scale empirical analyses.
Recent studies~\citep{geiping2025scaling,zhu2025scaling,du2025latent,mcleish2025teaching} have extended the recurrent Transformer architecture into language models. Huginn~\citep{geiping2025scaling} trained a 3.5B model from scratch, while Ouro~\citep{zhu2025scaling} introduced a more capable LoopLM, enabling it to compete with mainstream open-source LLMs. As a typical latent reasoning approach, LoopLM is expected to demonstrate test-time scaling. However, we find that rather than producing progressive gains with increased computation, performance typically peaks at a specific iteration depth and declines sharply beyond it. This phenomenon is especially pronounced in Ouro models.
% In contrast, Huginn’s model, which employs random loops during training, shows no severe degradation. 
% Current works lack in-depth analysis and insight into the training design, especially regarding the latent dynamics. Although Deep Equilibrium Models~\citep{bai2019deep,bai2021stabilizing} previously described the dynamics of looped models, they are not typical language model architectures and their training algorithms cannot be directly transferred. 
% Therefore, it is crucial to study the latent dynamics of current LoopLMs and address their performance degradation when increasing iterations during test time.
Existing studies lack a thorough analysis of LoopLM training design, particularly regarding latent dynamics. While DEQ~\citep{bai2019deep,bai2021stabilizing} examined the dynamics of looped networks, their architectures and training algorithms differ from standard language models. Therefore, investigating the latent dynamics in modern LoopLMs is crucial to address their unreliable scaling behavior.
